\subsection{Summary and outlook}

In this section, we have mainly presented results regarding the so-called inclusive radiative corrections.
These are suitable to estimate the effect of QED corrections on the differential decay widths (and consequently on the extraction of TFFs or their parameters), assuming that the emission of additional photons has been ignored during the analysis.
Any additional assumptions or cuts that are made in the experiment should be reflected in the calculation, so a close experiment--theory collaboration is essential for preparing the suitable tailored approach.

As pointed out earlier, the most desirable way to implement radiative corrections is as early as at the Monte Carlo level.
Already at that stage, one might typically observe that there are events missing in the MC, and the agreement with data is not satisfactory.
This leads to larger uncertainties and extracted parameters containing unsubtracted QED parts.

The corrections in Figs.~\ref{fig:Dalitz_RC} and \ref{fig:etap_e_RC} may seem large at first.
But one needs to realize that the Dalitz decays are dominated by a one-photon exchange (i.e., by the low-$x$ region due to a pole at vanishing photon virtuality).
Integrating over the Dalitz plot then leads to the fact that the QED corrections at the dilepton production threshold will be receiving the largest weights leading to the determination of the corrections to the integrated decay width that are $\sim$\;1\% as one would expect.
These resulting values can then be independently checked by employing different methods (see, e.g., Ref.~\cite{Lautrup:1971ew}), and it was done so for $\pi^0$ and $\Sigma^0$ Dalitz decays arriving at an exact agreement.
Note that using the soft-photon approximation, the corrections are typically negative all over the Dalitz plot, and the effects of the hard-photon corrections are instrumental for obtaining the observed positive total corrections (see, e.g., Ref.~\cite{Sidhu:1971myo}).

We have stated examples of processes for which the radiative corrections have been worked out or used in the experiment.
From the mentioned examples, it is straightforward to see the importance of radiative corrections and that with improving sensitivity and higher statistics of upcoming experimental set-ups, these will become increasingly indispensable in order to extract meaningful results from data.
The list of light-hadron decays for which the exact NLO QED radiative corrections are available can thus be expected to get longer.